\clearpage
\appendix
\titleformat{\section}
{\centering\normalfont\Large\bfseries}{\Alph{section}.}{.5em}{}
\renewcommand\thesubsection{\Alph{section}.\arabic{subsection}}
\section{Reproducibility and Evidence Inventory}

\subsection{Frozen models and adopted matrices}

The manifests in \path{handoff/v1_2/} identify the Q1-only and Q2/Q3
scoring artifacts.  The adopted Q2/Q3 matrices, raw samples, summaries,
and verification records are archived in
\path{handoff/Q123_v1/Q123_v1_integrated.zip}.  The paper's public tables
are reproduced from \path{q23_experiment_table.csv}, while independent
energy checks are recorded in \path{sdk_verification.json}.  The source-code
entry points and their respective roles are documented in
\path{appendix/source_code_inventory.md}.

\begin{table}[!htbp]
\centering
\caption{Selection-involved retrospective model comparison on the 3,969 labeled rows of file~03.}
\label{tab:q1-validation}
\small
\begin{tabularx}{\textwidth}{@{}Xrrrrrr@{}}
\toprule
Model & MAE & RMSE & Spearman & $\mathrm{MAE}_{V<5}$ & Top-25 & Top-50 \\
\midrule
Direct HGB                & 2.7078 & 4.6876 & 0.9267 & 11.2603 & 4 & 13 \\
Tail-weighted direct HGB  & 2.8655 & 4.7992 & 0.9223 &  8.3598 & 10 & 18 \\
Mean-gap property HGB     & 2.5278 & 4.6885 & 0.9377 &  6.3517 & 7 & 19 \\
Penalty-mixture HGB       & 2.4715 & 4.5827 & 0.9383 &  8.2654 & 8 & 20 \\
\bottomrule
\end{tabularx}
\end{table}

\begin{table}[!htbp]
\centering
\caption{Retrospective finite-budget replay on selection-involved file~03.
Random minima are exact distributional medians without replacement.}
\label{tab:q1-budget}
\small
\begin{tabularx}{\textwidth}{@{}r*{5}{>{\centering\arraybackslash}X}@{}}
\toprule
$K$ & Proxy $V<5$ hits & Random mean hits & Enrichment & Proxy min. $V$ & Random median min. $V$ \\
\midrule
10  &  6 & 0.1764 & 34.02 & 0.0196 & 12.6431 \\
25  & 13 & 0.4409 & 29.48 & 0.0196 & 10.9103 \\
50  & 23 & 0.8818 & 26.08 & 0.0196 &  3.4175 \\
100 & 34 & 1.7637 & 19.28 & 0.0196 &  1.2500 \\
\bottomrule
\end{tabularx}
\end{table}

\subsection{Precision comparisons and negative tests}

The historical Q2/Q3 paper-and-figures archive in \path{handoff/}
contains the floating-point 104- and 208-variable matrices, legal-state
energy checks, direct eight-bit quantization reports, two short-list
SPQC logs, and their decoding scripts.  The follow-up records in
\path{handoff/q23_fix/} correct the Ising constant by restoring the QUBO
offset once, compare full-pool and short-list penalties, and retain the
negative-input and regression tests.  The principal machine-readable
records supporting Sections~\ref{sec:q3-qubo}--
\ref{sec:results-discussion} are:

\begin{itemize}[leftmargin=2em]
 \item Stage 1 SDK--Ising record: QUBO--Ising round-trip,
       reference-spin dimensions, offsets, and energy residuals;
 \item Stage 2 quantization record: direct quantization
       failures and the three adopted 104-variable partitions;
 \item Stage 3 hit-rate record: short-list size,
       penalty, and local annealing-budget controls;
 \item Stage 8 negative-validation and Stage 9 regression records: rejected corrupt inputs and
       failure-path checks.
\end{itemize}

The negative-input audit rejected all 13 injected defect classes for
both Q2 and Q3.  Regression tests R1--R4 cover illegal-only solver output,
a missing SDK licence, complete Q2/Q3 candidate domains, and source-field
validation.  The mapping recheck used Kaiwu~1.4.4 without a solve licence;
the recorded CPU solves used Kaiwu SDK~1.4.3.  The floating-point Q2 and
joint-Q3 control QUBOs have SHA prefixes
\texttt{967b13d1957b} and \texttt{f3b447a1edb6}.  Their full hashes appear
in \path{appendix/source_code_inventory.md} and the raw manifests.  The
reported floating-point residuals do not apply to the adopted
104-variable integer matrices.

\begin{table}[!htbp]
\centering
\caption{Aligned scope of the three Q3 encodings.}
\label{tab:q3-encoding-comparison}
\scriptsize
\begin{tabularx}{\textwidth}{@{}lXXXX@{}}
\toprule
Encoding & Objective and oracle dependence & Variables and range & Ranking and legal coverage & Solves and role \\
\midrule
107 projection & Global pairwise fit; no known-winner input & 107; $\max|h|=254$ & Approximate; 8,112 states, Top-5 overlap 0 & One; adverse low-tail control \\
208 joint & Exact floating proxy cost; no known-winner input & 208; $\max|h|=161.187$, $\max|J|=1$ & Exact before quantization; 8,112 states & One; floating precision control \\
104 groupwise & Two-level cost uses known group argmin & 104/group; $\max|h|=102$, $\max|J|=1$ & Retains argmin only; 2,704 states/group & Three plus raw-score comparison; adopted mapping \\
\bottomrule
\end{tabularx}
\end{table}

\begin{table}[!htbp]
\centering
\caption{Unified mapping audit.  SHA prefixes identify the audited QUBO.}
\label{tab:mapping-audit}
\scriptsize
\begin{tabularx}{\textwidth}{@{}lXXXXX@{}}
\toprule
Object & Version & Matrix & Comparison & Offset & Maximum discrepancy / scope \\
\midrule
Adopted 104 integer Q2/Q3 & SDK 1.4.3 archive & final hashes below & formula--matrix, all legal states & $+4$ & 0 integer units; exact discrete check \\
Q2 floating control & Kaiwu 1.4.4 & \texttt{967b13d1} & SDK coefficient / energy chain & $+4$ & $2.13\times10^{-14}$ / $1.42\times10^{-12}$; numerical tolerance \\
Q3 joint 208 control & Kaiwu 1.4.4 & \texttt{f3b447a1} & SDK coefficient / energy chain & $+4$ & $9.95\times10^{-14}$ / $4.57\times10^{-12}$; numerical tolerance \\
Adopted 104 under 1.4.4 & NA & NA & comparable floating residual & NA & NA; 208 values not substituted \\
\bottomrule
\end{tabularx}
\end{table}

In Table~\ref{tab:mapping-audit}, the first residual in each floating-point
row is $\epsilon_{\mathrm{coef}}$, the maximum absolute difference between
SDK-reconstructed and independently derived Ising coefficients; the second
is $\epsilon_{\mathrm{energy}}$, the maximum offset-corrected energy-chain
residual over the audited states.  The adopted integer matrices instead
have zero formula-to-stored-matrix discrepancy in integer energy units.
These are distinct from the ranking and tie metrics used to assess
coefficient quantization.

\subsection{Projection and blind-pool evidence}

The 107-variable projection matrix and its 300 local-SA samples are
archived under \path{handoff/q23_supplement/O10/}.  Its optimum is
Rb--Ta--O3, whereas the original proxy optimum is Na--Nb--O3; the latter
appears in only 2 of the 300 samples.  This experiment is a projection
control, not an optimum-preservation result.

The dedicated file~11/12 records are stored in
\path{handoff/q23_supplement/O14/}.  Each floating-point and integer model
includes 300 samples with both bit and spin representations.  The matrix
SHA-256 values are:

\begin{itemize}[leftmargin=2em]
 \item Q2/NO2 float: \texttt{eeef8d05f7d4...21938311};\par
       integer: \texttt{5ba756450c57...bd797a2}.
 \item Q3/O3 float: \texttt{9749bb7e5180...4d6f346};
       integer: \texttt{aee1caab0b86...671af49}.
 \item Q3/NO2 uses the same float/integer hashes as Q2/NO2 because the
       candidate-score table and construction coincide for that group.
 \item Q3/FO2 float: \texttt{aea073fccfa8...74f6ee21};\par
       integer: \texttt{50ae41a234c7...ab359f1}.
\end{itemize}

The floating-point/eight-bit batch-hit counts are respectively
25/16 (Q2 NO2), 27/13 (Q3 O3), 25/16 (Q3 NO2), and 26/11 (Q3 FO2);
each batch is one seed with ten reads, not a single-sample hit rate.
All eight blind-pool models returned 300/300 legal samples.  The
corresponding eight-bit minimum-tier tie counts are 11, 19, 11, and 34
for Q2 NO2, Q3 O3, Q3 NO2, and Q3 FO2, respectively; each floating-point
minimum was unique.  The audited local simulated-annealing call times
range from 38.532520 to 48.535693~s across the eight runs.  These times
cover the local annealing call only, not matrix construction, conversion,
decoding, or feasibility checks.  The older 45--786~s range was withdrawn
because no corresponding local run record was found.

\subsection{Final SPQC task evidence and unresolved fields}

The four final 104-variable task receipts are archived under
\path{handoff/q23_supplement/O13/}:

\begin{itemize}[leftmargin=2em]
 \item Q2/NO2 task \texttt{c249dca6...f4ebf2a}, matrix
       \texttt{fee2320be719...81bee5c};
 \item Q3/O3 task \texttt{27a3df46...dc1a59f}, matrix
       \texttt{26d234efa34c...a8a82ef};
 \item Q3/NO2 task \texttt{1e00e1fb...c9477d2}, matrix
       \texttt{fee2320be719...81bee5c};
 \item Q3/FO2 task \texttt{c528ea69...e99ee8d}, matrix
       \texttt{5b49510e7821...6fb7309}.
\end{itemize}

\begin{table}[!htbp]
\centering
\caption{Final 104-variable SPQC-550 task returns.  Total denotes the
creation-to-finish interval, not isolated execution time.}
\label{tab:spqc-final-results}
\small
\begin{tabularx}{\textwidth}{@{}lXrrrr@{}}
\toprule
Task & Group & Returned & Legal & Argmin hit & Total / evolution \\
\midrule
Q2 & NO2 & 10 & 1 & yes & 22 s / 15.0116 ms \\
Q3 & O3  & 10 & 4 & yes & 74 s / 9.7488 ms \\
Q3 & NO2 & 7  & 0 & no  & 112 s / 0.6991 ms \\
Q3 & FO2 & 5  & 1 & yes & 16 s / 7.3587 ms \\
\bottomrule
\end{tabularx}
\end{table}

Historical tasks \texttt{q2\_int\_01} and \texttt{q3\_int\_01} used 17
and 30 variables, respectively, and are kept separate from the final
104-variable submissions.  The credit ledger records seven debits, with
the balance falling from 31 to 24, while six task receipts are archived.
The \texttt{probe4\_01} entry has no original receipt, so its task details
remain NA.  Platform queue time, complete execution time, monetary
billing, and physical labels for all proxy winners also remain NA.

The remaining provenance qualifications are recorded without
overstating the evidence.  The historical stage-2 0/300 count is associated
with a demonstrated column-index bug: on the same 300 samples, the
corrected counter gives 8/300 while the legacy counter gives 0/300.
Because the original historical run log is absent, that historical value
remains \textbf{UNVERIFIED}; the new O11-R01/R02/R03 reconciliations are
reported separately.  The earlier 162.556~s local-SA total has no source
package in the repository and remains \textbf{UNVERIFIED}; the audited
O15-R01--R04 total is 178.131869~s.  The O-14 45--786~s range was withdrawn
and replaced by the eight recorded annealing-call times of
38.532520--48.535693~s.  Finally, 0.48--0.69~s refers only to 30
\texttt{opt.solve()} calls and excludes setup, conversion, adaptation,
decoding, and feasibility checks; it is not comparable to the separately
reported 20.81--20.96~s full-workflow SDK time.  Historical 20.04~ms
records belong to a different configuration.

\subsection{Pipeline and resource-cost audit}
\label{sec:appendix-cost-audit}

The archived intake tests injected 13 defect classes: external
identifiers; incorrect compositions, pool fields, or source-pool fields;
an incorrect source field alone; missing pool or source columns; duplicate
or missing rows; NaN or negative scores; mixed model versions; and
constant prediction tables.  Both Q2 and Q3 rejected all 13 classes.
Regression test R1 returns \texttt{solver\_failed} with no winner when all
samples are illegal.  R2 detects a missing SDK licence.  R3 preserves all
2,704 Q2 candidates and all 2,704 candidates in each Q3 group through the
joint entry point.  R4 rejects invalid source fields.

The measured complete QUBO build--convert--adapt--check pipeline
includes matrix construction and took 6.551~ms for Q2, 7.430, 7.524,
and 7.591~ms for the three Q3 groups, and 23.197~ms for the separately
measured three-group Q3 run.  Each single-group call constructed one
matrix; the combined Q3 call constructed three.  Kaiwu SDK~1.4.3 CPU
full-workflow solves took 20.81--20.96~s per group (83.413~s in total)
in the external review environment.  A separate audit measured
0.48--0.69~s for 30 \texttt{opt.solve()} calls only; this excludes
initialization, conversion, adaptation, decoding, and feasibility checks
and is not comparable to the full-workflow timing.  The audited local-SA
runs O15-R01--R04 took 50.289327, 41.845110, 40.499950, and
45.497482~s, totaling 178.131869~s.  The earlier 162.556~s total remains
\textbf{UNVERIFIED} and is excluded.  The SPQC-550 creation-to-finish
intervals were 22, 74, 112, and 16~s; the reported single-evolution times
were 15.0116, 9.7488, 0.6991, and 7.3587~ms.  Neither field isolates
queueing or complete execution.  Monetary billing, queue time, complete
execution time, and a comparable CPU-to-device end-to-end cost remain NA.
